\chapter{Upravljanje raspodijeljenim tijekovima}
\label{ch:tijekovi}

Poglavlje~\ref{ch:raspodijeljeni} govorilo je o jednom pozivu. Ovo poglavlje
govori o nizu koraka koji zajedno čine jedan posao. Utvrđuje kako se stanje
drži razumljivim kada je posao djelomično gotov, a nastavak neizvjestan.

% ------------------------------------------------------------
\section{Granice klasičnih transakcija}
\label{sec:granica-transakcije}

Transakcija je pojam stariji od računarstva i Gray ga tako i uvodi. Njegov
primjer je ugovor: pregovori traju, a onda ih jedan čin
zaključi \cite{gray1981transaction}. Iz te analogije izvodi tri svojstva.
Konzistentnost znači da transakcija poštuje dopuštena pravila. Atomarnost
znači da se dogodi ili se ne dogodi. Trajnost znači da se jednom počinjena
transakcija ne može poništiti.

U tom popisu izolacije nema. Četveroslovni
akronim ACID, koji se danas navodi kao da je oduvijek postojao, uveli su
dvije godine kasnije Härder i Reuter \cite{haerder1983principles}.

Gray odmah dodaje ogradu koja je za ovaj rad važnija od samih svojstava.
Ugovor se može poništiti samo ako je od početka bio nezakonit. Inače se
loša transakcija ispravlja daljnjim transakcijama.

Prijelaz s jednog sudionika na više njih rješava dvofazno počinjanje, koje
Gray izlaže u istoj bilježnici u kojoj i paradoks dvaju
generala \cite{gray1978notes}. Koordinator u prvoj fazi pita sve sudionike
jesu li spremni. Sudionik koji pristane mora prije odgovora trajno zapisati
dovoljno podataka da poslije može i poništiti i ponoviti. U drugoj fazi
koordinator zapisuje odluku i šalje je svima.

Uspjeh tog protokola Gray pripisuje jednoj stvari: odluka o počinjanju
središnja je i nije vremenski ograničena. Ta rečenica sadrži i cijenu.
Sudionik koji je glasao za počinjanje čeka presudu i sam je ne smije
donijeti. Ako koordinator u tom trenutku otkaže, sudionik po Grayevu
opisu pri oporavku pita koordinatora je li transakcija počinjena ili
prekinuta. Dok koordinatora nema, sudionik ne zna, a resursi koje drži
ostaju zaključani.

Bailis i Ghodsi to izriču izravno za razdvojenu mrežu. Kada dva procesa ne
mogu komunicirati, izmjene se ne mogu sinkrono proširiti na sve, a da se
pritom ne
blokira \cite{bailis2013eventual}.

Za studiju slučaja iz poglavlja~\ref{ch:studija} rasprava završava ranije
nego što bi trebala. Dvofazno počinjanje preko dva kanala nije isključeno
zato što je skupo ili zato što bi koordinator blokirao. Isključeno je zato
što nijedno od dvaju sučelja u takvom protokolu ne sudjeluje. Ni razmjena
preko pristupnih točaka ni fiskalizacijski poziv ne nude fazu pripreme ni
glasanje, pa presudu koordinatora nema kome uputiti.

Propis to tehničko stanje dodatno omeđuje. Članak 48.\ stavak 3.\ nalaže da
se fiskalizacija provodi odvojeno od postupka razmjene, dakle kao dva
postupka, a ne kao jedan. Rad iz te odredbe ne izvodi zabranu svake
koordinacije dvaju postupaka u lokalnoj bazi, jer odredba o transakcijskoj
organizaciji izvedbe ne govori. Izvodi uže: zajedničkog protokola atomarnog
dovršetka nema, pa izvedba mora računati na međustanja u kojima je jedan
postupak dovršen, a drugi nije.

Ono što preostaje kada atomarnost otpadne obično se zove eventualnom
konzistentnošću. Bailis i Ghodsi je definiraju kao jamstvo da će, ako nema
novih izmjena nekog podatka, sva čitanja tog podatka na kraju vratiti istu
vrijednost. Isti autori ne skrivaju koliko je to slabo. Ni u jednom
trenutku korisnik ne može isključiti nekonzistentno ponašanje, a jedino
jamstvo je da će se u nekom trenutku u budućnosti dogoditi nešto dobro.

Slabo jamstvo ipak nije jedino dostupno. Autori pokazuju da se uz visoku
raspoloživost mogu zadržati i jača svojstva, među njima uzročna
konzistentnost. Navode i rezultat po kojem je ona gornja granica onoga što se
pod razdvajanjem mreže može postići. Za ovaj rad iz
toga slijedi metodološko pravilo. Reći da sustav nije konzistentan ne znači
mnogo; treba imenovati koje jamstvo daje.

% ------------------------------------------------------------
\section{Saga i kompenzacija}
\label{sec:saga}

Kada atomarnost preko granice otpadne, popravljanje prelazi iz protokola u
poslovnu logiku. Gray taj prijelaz imenuje \cite{gray1981transaction}:

\begin{quote}
„Nakon što se transakcija počini, njezini se učinci mogu izmijeniti samo
izvođenjem daljnjih transakcija. Na primjer, ako je netko premalo plaćen,
ispravak je izvesti drugu transakciju koja isplaćuje dodatni iznos. Takve
naknadne transakcije zovu se kompenzacijskim
transakcijama.“\footnote{„\emph{Once a transaction commits, its effects can
only be altered by running further transactions. For example, if someone is
underpaid, the corrective action is to run another transaction which pays an
additional sum. Such post facto transactions are called compensating
transactions.}“ \cite{gray1981transaction}} \end{quote}

Primjer kojim Gray pokazuje granice klasične transakcije je putnička
agencija. Agent dogovara let, pa automobil, pa hotel, pa naplaćuje karticu.
Kupac zatim otkaže putovanje. Svaka pojedina interakcija bila je transakcija
s tuđom organizacijom. Tu Gray izvodi rečenicu koja opisuje svaki tijek preko
vanjskog sudionika. Agent ne može jednostrano prekinuti interakciju nakon što
je dovršena, nego mora izvesti kompenzacijsku transakciju koja prethodnu
poništava.

Garcia-Molina i Salem taj obrazac su formalizirali pod imenom
saga \cite{garciamolina1987sagas}. Dugotrajna transakcija je saga ako se može
napisati kao niz transakcija koje se smiju izmjenjivati s drugima. Sustav
jamči da će ili sve transakcije u sagi uspjeti, ili će se izvesti
kompenzacijske transakcije koje djelomično izvođenje isprave. Cijena je
poznata: međustanja sage vide se izvana, pa saga nije izolirana.

Izvedba sage traži i odluku o tome tko vodi niz. Orkestracija stavlja jednog
voditelja koji poziva korake redom. Koreografija razdjeljuje odluku na
sudionike koji reagiraju na
događaje \cite{megargel2021orchestration}. Za tijek koji ovaj rad opisuje
izbor slijedi iz propisa, ne iz ukusa. Za ispravno izdavanje, zaprimanje i
fiskalizaciju odgovara porezni obveznik (čl.~58.~st.~2.). Kada odgovornost
za ishod nosi jedan sudionik, niz vodi taj sudionik.

Poglavlje~\ref{ch:studija} ovisi o tome ima li u fiskalizaciji što
kompenzirati.

Odgovor daje Grayeva podjela sastavnica transakcije na četiri vrste:

\begin{itemize}[nosep, leftmargin=1.4cm]
  \item radnje nad nezaštićenim objektima,
  \item zaštićene radnje, koje se mogu poništiti i ponoviti,
  \item ugniježđene transakcije, koje se poništavaju kompenzacijom,
  \item stvarne radnje, za koje kaže da se mogu odgoditi, ali ne i poništiti.
\end{itemize}

Fiskalizacija pripada posljednjoj vrsti. Dostava podataka Poreznoj upravi
zakonska je obveza, ne rezervacija. Propis ne predviđa poruku kojom bi obveznik
povukao provedenu fiskalizaciju. Predviđa odgodu, i to omeđenu: pet radnih
dana kod tehničke nemogućnosti (čl.~49.~st.~1.).

Na fiskalizaciju je primjenjiv onaj dio obrasca sage koji prati niz lokalnih
koraka i evidentira njihov napredak. Kompenzacija nije primjenjiva, jer koraci
nisu poništivi. Poglavlje~\ref{ch:studija} iz obrasca preuzima odgodu uz zapis
o napretku.

Helland i Campbell isti zaključak izvode iz druge strane. Kada se poslovno
pravilo provodi asinkrono, ono postaje vjerojatnosno, jer sustav koji ga
provodi ne zna hoće li biti dostupan da ga
provede \cite{helland2009quicksand}. Njihov zaključak je da sustav ili snosi
trošak sinkrone potvrde, ili se povremeno mora ispričati. Fiskalizaciji drugi
izlaz nije dostupan. Propis ne predviđa naknadnu ispriku kao način rješavanja
propuštene prijave, a propuštanje roka je prekršaj (čl.~73.~st.~1.~t.~21.).

Iz toga slijedi granica koju odjeljak~\ref{sec:dizajn} razrađuje. Skup
poslovnih poruka ne sadrži nijednu koja bi provedenu fiskalizaciju poništila,
pa kompenzacija nema predmet.

% ------------------------------------------------------------
\section{Stanje kao automat}
\label{sec:automat}

Ako se posao ne može zatvoriti u jednu transakciju, mora se moći opisati
kao stanje. Schneider daje karakterizaciju koja je za to
dovoljna \cite{schneider1990statemachine}:

\begin{quote}
„Izlazi automata stanja potpuno su određeni nizom zahtjeva koje obrađuje,
neovisno o vremenu i o svakoj drugoj aktivnosti u
sustavu.“\footnote{„\emph{Outputs of a state machine are completely
determined by the sequence of requests it processes, independent of time and
any other activity in a system.}“ \cite{schneider1990statemachine}}
\end{quote}

Dvije riječi u toj rečenici nose posljedice. Prva je niz. Stanje ovisi o
redoslijedu obrađenih zahtjeva, pa se dva sustava koja su obradila isti niz
nalaze u istom stanju.

Odatle se lako izvede više nego što u rečenici stoji. Determinizam nije
idempotentnost. Automat koji vodi brojač i dvaput primi istu naredbu povećat
će brojač dvaput, jer su to dva zahtjeva u nizu, a ne jedan. Kako ponavljanje
ne bi mijenjalo stanje, poruka mora nositi identifikator zahtjeva, a automat
mora pamtiti koje je zahtjeve već obradio.

Schneider to i razdvaja: jedno je da niz zahtjeva određuje izlaze, a drugo da
svi primjerci obrade zahtjeve u istom redoslijedu. Popuštanje uvjeta redoslijeda
Schneider veže uz komutativnost operacija. Jednak skup poruka u drukčijem
redoslijedu zato ne daje jednako stanje.

Druga je vrijeme. Izlazi ne ovise o vremenu. Schneider uz uvjete redoslijeda
dodaje napomenu da zbog kašnjenja u mreži automat neće nužno obraditi
zahtjeve redom kojim su poslani ni redom kojim su primljeni.

Tu se model razilazi s propisom, a poglavlje~\ref{ch:studija} tu razliku
nasljeđuje. Model je određen nizom poruka, a rokovi u propisu mjere se satom. „U trenutku
izdavanja“ i „pet radnih dana“ su vremenske veličine. Model automata ih ne
zabranjuje. Protek roka može se modelirati kao izričit događaj, koji netko
mora dostaviti, i tada je redovan član niza. Može se modelirati i kao uvjet
nad prijelazom, čime prijelaze pokreću samo poruke, a rok odlučuje je li
prijelaz dopušten. Poglavlje~\ref{ch:implementacija} koristi oboje i
obrazlaže zašto, ali oba su izvedbena izbora, a ne jedino teorijski dopušteno
rješenje.

Helland ta dva svijeta spaja u
praksi \cite{helland2007lifebeyond}. Sustav se dijeli na entitete, svaki sa
svojim opsegom serijalizacije, a među njima putuju poruke koje se šalju
barem jednom i obrađuju idempotentno. Transakcija postoji, ali samo unutar
entiteta. Preko granice entiteta postoje poruke i zapamćeno stanje o tome
što je već obrađeno.

Da bi se o stanjima uopće govorilo kao o uređenima, potreban je pojam
uzročnog prethođenja, koji je uveo Lamport \cite{lamport1978time}.

Odatle slijedi model koji poglavlje~\ref{ch:studija} treba.
Životni ciklus dokumenta opisuje se kao automat čija su stanja izrijekom
imenovana, uključujući ona koja propis ne imenuje. Stanje
„poslano, nepotvrđeno“ nije rupa u izvedbi. To je stanje u kojem se sustav
nalazi dok čeka odgovor koji možda nikada ne stigne.
